\title{Group Sequence Policy Optimization}

\begin{abstract}
We present Group Sequence Policy Optimization (GSPO), a robust and effective reinforcement learning method tailored for training large language models. Departing from earlier approaches that utilize token-level importance ratios, GSPO instead formulates the importance ratio over the entire sequence, enabling sequence-level operations such as clipping, rewarding, and optimization. Our empirical results indicate that GSPO surpasses the GRPO algorithm in both training efficiency and overall performance, significantly enhances stability in Mixture-of-Experts (MoE) RL training scenarios, and offers promising avenues for streamlining RL infrastructure design. The advantages of GSPO have played a key role in achieving the significant advancements observed in the latest Qwen3 models.
\end{abstract}

\section{Introduction}
\label{sec:intro}

Reinforcement learning (RL) has become a fundamental approach for advancing the performance of language models \citep{o1,dpsk-r1,qwq32b,qwen3}. When leveraged at scale, RL enables language models to perform advanced reasoning, equipping them to address complex tasks like competitive mathematics and programming by supporting more intricate and extended chains of reasoning.

Achieving further scaling of RL via increased computational resources crucially depends on the ability to preserve stable and resilient training processes. Nevertheless, even leading RL methods, such as GRPO \citep{grpo}, currently suffer from pronounced instability when deployed on extremely large language models. This often leads to abrupt and unrecoverable collapses in model performance \citep{qwen3, minimax-m1}. Such instability represents a significant obstacle for extending the limits of what language models can achieve through ongoing RL-based training.

In this study, we demonstrate that the inherent instability of GRPO can be traced to a foundational flaw: the improper use and breakdown of importance sampling weights within its algorithmic structure. This defect generates substantial variance in the training signal, which accumulates as responses become longer, and is exacerbated by the clipping strategy, culminating in model breakdown.

To overcome these fundamental issues, we introduce \textbf{Group Sequence Policy Optimization (GSPO)}, a novel reinforcement learning method tailored for large language model training. The principal advancement in GSPO is its rigorously derived importance ratio, calculated using sequence likelihood in accordance with established importance sampling theory \citep{click}. Moreover, GSPO determines normalized rewards by evaluating the advantages across multiple responses for each query, thereby maintaining coherence between the granularity of rewards and the policy optimization process.

Our experimental results reveal that GSPO consistently outperforms GRPO in terms of training stability, efficiency, and overall effectiveness. Notably, GSPO intrinsically addresses the stability issues commonly faced during RL training of extensive Mixture-of-Experts (MoE) architectures, thereby obviating the requirement for elaborate stabilization mechanisms and offering prospects for more streamlined RL systems. These advantages directly underpin the marked enhancements observed in the newest Qwen3 model series. We anticipate that GSPO will serve as a reliable and scalable cornerstone for future large-scale RL endeavors involving language models.

\section{Preliminaries}

\paragraph{Notation}
Throughout this work, we represent an autoregressive language model with parameters $\theta$ as a policy $\pi_\theta$. The symbol $x$ refers to a query, while $\mathcal{D}$ signifies the collection of all queries. For a given response $y$ to query $x$, its probability according to the policy $\pi_\theta$ is given by $\pi_\theta (y | x)=\prod_{t=1}^{|y|} \pi_\theta (y_t | x, y_{<t} )$, where $|y|$ represents the total number of tokens in the response $y$. Each pair of query and response $(x, y)$ is evaluated by a verifier $r$, yielding a reward value $r(x, y) \in [ 0, 1 ]$.

\paragraph{Proximal Policy Optimization (PPO)} Using samples generated from the old policy $\pi_{\theta_\text{old}}$, PPO \citep{ppo} constrains the policy update within a proximal region of the old policy through the clipping mechanism. Specifically, PPO employs the following objective for policy optimization (we omit the KL regularization term hereinafter for brevity, as it is not the focus of this paper):

\begin{align}
\mathcal{J}_\text{PPO}(\theta) = \mathbb{E}_{ x \sim \mathcal{D},\, y \sim \pi_{\theta_\text{old}}( \cdot | x) }
\left[ \frac{1}{|y|} \sum_{t=1}^{|y|} 
\min \left( w_{t}(\theta) \widehat{A}_{t},  \, \mathrm{clip} \left( w_{t}(\theta), 1 - {\varepsilon}, 1 + {\varepsilon}\right) \widehat{A}_{t} \right)
\right],
\end{align}

In this context, the token importance ratio $w_{t}(\theta)$ is given by $w_{t}(\theta) = \frac{ \pi_{\theta} (y_{t} | x, y_{<t}) }{ \pi_{\theta_\text{old}} (y_{t} | x,y_{<t})}$, the advantage $\widehat{A}_{t}$ associated with $y_t$ is derived from a separate value model, and $\varepsilon$ determines the allowable clipping range for these importance ratios.

A fundamental obstacle when deploying PPO stems from its significant dependence on the value model. Typically, the value model's architecture closely matches that of the policy model, resulting in increased demands on memory and computational resources. In addition, the success of PPO algorithms relies on accurate estimations produced by the value model. Attaining a trustworthy value model is already a nontrivial task; scaling such models to handle lengthier sequences and more sophisticated problems amplifies this difficulty.

\paragraph{Group Relative Policy Optimization (GRPO)} GRPO \citep{grpo} bypasses the need for the value model by computing the relative advantage of each response within a group of responses to the same query. Specifically, GRPO optimizes the following objective:

\begin{align}
\label{equ:grpo}
\mathcal{J}_\text{GRPO}(\theta) = \mathbb{E}_{ x \sim \mathcal{D},\, \{y_i\}_{i=1}^G \sim \pi_{\theta_\text{old}}( \cdot | x) }
\left[ \frac{1}{G} \sum_{i=1}^{G} \frac{1}{|y_i|} \sum_{t=1}^{|y_i|} 
\min \left( w_{i,t}(\theta) \widehat{A}_{i,t},  \, \mathrm{clip} \left( w_{i,t}(\theta), 1 - {\varepsilon}, 1 + {\varepsilon}\right) \widehat{A}_{i,t} \right)
\right],
\end{align}

where $G$ is the number of generated responses to each query $x$ (i.e., the group size), and the importance ratio $w_{i,t}(\theta)$ and advantage $\widehat{A}_{i,t}$ of token $y_{i,t}$ are:

\begin{align}
    w_{i,t}(\theta)=\frac{ \pi_{\theta} (y_{i,t} | x, y_{i,<t}) }{ \pi_{\theta_\text{old}} (y_{i,t} | x,y_{i,<t})},\quad\ 
    \widehat{A}_{i,t} = \widehat{A}_{i} = \frac{r(x, y_i) - \mathrm{mean} \left( \{ r(x, y_i) \}_{i=1}^G \right) }{ \mathrm{std} \left( \{ r(x, y_i) \}_{i=1}^G \right) },
\end{align}

respectively, where all the tokens in $y_i$ share the same advantage as $\widehat{A}_{i}$.

\section{Motivation}
\label{sec:motivation}

As models increase in size, exhibit greater sparsity (such as in Mixture-of-Experts architectures), and generate longer responses, larger rollout batch sizes become essential for achieving optimal hardware throughput during reinforcement learning. To make better use of collected samples, it is common to divide these large rollout batches into several mini-batches when performing gradient computations. However, this strategy inherently shifts the setting to off-policy learning, since response samples $y$ originate from a previous policy $\pi_{\theta_\text{old}}$ and not the current policy $\pi_\theta$ under training. This circumstance motivates the need for the clipping strategies used in PPO and GRPO, which serve to exclude samples that deviate excessively from the current policy and might otherwise adversely affect gradient estimation.

While mechanisms like clipping aim to manage this off-policy discrepancy, we identify a more fundamental issue in GRPO: \textit{its objective is ill-posed}. This problem becomes particularly acute when training large models on long-response tasks, leading to catastrophic model collapse. The ill-posed nature of the GRPO objective stems from a misapplication of importance sampling weights. The principle of importance sampling is to estimate the expectation of a function $f$ under a target distribution $\pi_\text{tar}$ by re-weighting samples drawn from a behavior distribution $\pi_\text{beh}$:

\begin{align}
\mathbb{E}_{z \sim \pi_\text{tar}} \left[ f(z) \right] = \mathbb{E}_{z \sim \pi_\text{beh}} \left[ \frac{ \pi_\text{tar} (z) }{ \pi_\text{beh} (z)} \, f(z) \right].
\end{align}

A key aspect here is that successful distribution correction depends on averaging across many samples ($N \gg 1$) drawn from the behavior policy $\pi_\text{beh}$; under these conditions, the importance weight $\frac{ \pi_\text{tar} (z) }{ \pi_\text{beh} (z)}$ can compensate for the difference between the sampling and target distributions.

GRPO, by contrast, implements the importance weight $\frac{ \pi_{\theta} (y_{i,t} | x, y_{i,<t}) }{ \pi_{\theta_\text{old}} (y_{i,t} | x,y_{i,<t})}$ at every token $t$. However, this ratio is calculated using only a single instance $y_{i,t}$ for each subsequent-token probability $\pi_{\theta_\text{old}}( \cdot | x, y_{i, <t})$, so it fails to accomplish proper correction between behavior and target policies. The result is an introduction of substantial high-variance noise to the gradient estimates during training; this effect compounds throughout lengthy sequences and is further intensified by the use of clipping. Empirically, we have witnessed that this often precipitates a collapse of the model, after which recovery becomes infeasible. Subsequent efforts—such as reloading a prior checkpoint, extensive hyperparameter searches (including adjusting clipping values), increasing the length of generations, or changing the RL tasks—have not succeeded in revitalizing training.

These findings underscore a structural flaw in GRPO’s methodology. The breakdown of the token-wise importance correction highlights a fundamental insight: \textbf{the scope of the optimization objective should coincide with the scope of the reward}. Because rewards are assigned for the full sequence, applying correction at each token is inherently ill-suited. Accordingly, we are compelled to abandon token-level optimization and instead investigate strategies for applying importance weighting and conducting optimization at the \textit{sequence level}.

\section{Algorithm}

\subsection{GSPO: Group Sequence Policy Optimization}

Although the token-level importance weight $\frac{ \pi_{\theta} (y_{i,t} | x, y_{i,<t}) }{ \pi_{\theta_\text{old}} (y_{i,t} | x,y_{i,<t})}$ presents issues in GRPO, it is worth noting that, for language generation tasks, the \textit{sequence-level} importance weight $\frac{ \pi_{\theta} (y | x) }{ \pi_{\theta_\text{old}} (y | x)}$ carries an interpretable theoretical significance. Specifically, it measures the extent to which a response $y$, generated from $\pi_{\theta_\text{old}} (\cdot | x)$, diverges from what would be generated by $\pi_{\theta} (\cdot | x)$; this connection is well-aligned with sequence-level rewards, and additionally, it provides meaningful guidance for the application of the clipping procedure.

Based on this straightforward observation, we propose the \textbf{Group Sequence Policy Optimization (GSPO)} algorithm. GSPO employs the following sequence-level optimization objective:

\begin{align}
\mathcal{J}_\text{GSPO} (\theta) =
\mathbb{E}_{ x \sim \mathcal{D},\, \{y_i\}_{i=1}^G \sim \pi_{\theta_\text{old}}( \cdot | x) }
\left[ 
\frac{1}{G} \sum_{i=1}^{G}
\min \left( s_{i}(\theta)  \widehat{A}_{i},  \, \mathrm{clip} \left( s_{i}(\theta), 1 - {\varepsilon}, 1 + {\varepsilon} \right) \widehat{A}_{i} \right) 
\right],
\label{equ:gspo}
\end{align}

where we adopt the group-based advantage estimation:

\begin{align}
\widehat{A}_{i} = \frac{r(x, y_i) - \mathrm{mean} \left( \{ r(x, y_i) \}_{i=1}^G \right) }{ \mathrm{std} \left( \{ r(x, y_i) \}_{i=1}^G \right) },
\end{align}

and define the importance ratio $s_{i}(\theta)$ based on sequence likelihood \citep{click}:

\begin{align}
s_{i}(\theta) = \left( \frac{ \pi_{\theta} (y_i | x) }{ \pi_{\theta_\text{old}} (y_i | x)} \right)^{\frac{1}{|y_i|}}
=
\exp \left( \frac{1}{|y_i|} \sum_{t=1}^{|y_i|} \log \frac{ \pi_{\theta} (y_{i,t} | x, y_{i,<t}) }{ \pi_{\theta_\text{old}} (y_{i,t} | x,y_{i,<t})} \right).
\end{align}

As a result, GSPO conducts clipping at the full sequence level rather than per token, thereby filtering out highly ``off-policy'' samples during gradient estimation; this approach is coherent with both sequence-level rewards and the corresponding optimization framework. Importantly, we employ length normalization for $s_{i}(\theta)$ to mitigate variance and to keep $s_{i}(\theta)$ within a consistent numeric interval. If omitted, substantial likelihood shifts in just a handful of tokens could induce significant volatility in the sequence-level importance ratio, necessitating different clipping intervals for responses of varying lengths. Additionally, we observe that the typical magnitude of clipping ranges in GSPO, when compared to those used in earlier methods such as GRPO, differs considerably, which arises from the distinct ways in which importance ratios are formulated in these algorithms.

\subsection{Gradient Analysis}
\label{subsec:gradient}

We can derive the gradient of the GSPO objective as follows (clipping is omitted for brevity):

\begin{align}
\nabla_{\theta} \mathcal{J}_\text{GSPO} (\theta)
=&\ 
\nabla_{\theta} \mathbb{E}_{ x \sim \mathcal{D},\, \{y_i\}_{i=1}^G \sim \pi_{\theta_\text{old}}( \cdot | x) }
\left[ 
\frac{1}{G} \sum_{i=1}^{G} s_{i}(\theta) \widehat{A}_{i}
\right] \\
= &\ 
\mathbb{E}_{ x \sim \mathcal{D},\, \{y_i\}_{i=1}^G \sim \pi_{\theta_\text{old}}( \cdot | x) }
\left[ 
\frac{1}{G} \sum_{i=1}^{G}
s_{i}(\theta) \widehat{A}_{i}
\cdot \nabla_{\theta} \log s_{i}(\theta)
\right] \\
=&\ 
\mathbb{E}_{ x \sim \mathcal{D},\, \{y_i\}_{i=1}^G \sim \pi_{\theta_\text{old}}( \cdot | x) }
\left[ 
\frac{1}{G} \sum_{i=1}^{G} 
\left( \frac{ \pi_{\theta} (y_i | x) }{ \pi_{\theta_\text{old}} (y_i | x)} \right)^{\frac{1}{|y_i|}} \widehat{A}_{i} 
\cdot \frac{1}{|y_i|} \sum_{t=1}^{|y_i|} \nabla_{\theta} \log \pi_{\theta} (y_{i,t} | x, y_{i,<t}) 
\right].
\label{equ:gspo_grad}
\end{align}

For comparison, the gradient of the GRPO objective is as follows (note that $\widehat{A}_{i,t} = \widehat{A}_{i}$):

\begin{align}
\nabla_{\theta} \mathcal{J}_\text{GRPO} (\theta) 
=&\ 
\nabla_{\theta} \mathbb{E}_{ x \sim \mathcal{D},\, \{y_i\}_{i=1}^G \sim \pi_{\theta_\text{old}}( \cdot | x) }
\left[ \frac{1}{G} \sum_{i=1}^{G} \frac{1}{|y_i|} \sum_{t=1}^{|y_i|} 
w_{i,t}(\theta) \widehat{A}_{i,t}
\right] \\
=&\
\mathbb{E}_{ x \sim \mathcal{D},\, \{y_i\}_{i=1}^G \sim \pi_{\theta_\text{old}}( \cdot | x) }
\left[ \frac{1}{G} \sum_{i=1}^{G} \widehat{A}_{i} 
\cdot \frac{1}{|y_i|} \sum_{t=1}^{|y_i|} 
\frac{ \pi_{\theta} (y_{i,t} | x, y_{i,<t}) }{ \pi_{\theta_\text{old}} (y_{i,t} | x,y_{i,<t})} 
\nabla_{\theta} \log \pi_{\theta} (y_{i,t} | x, y_{i,<t})  
\right].
\end{align}

The principal difference between GSPO and GRPO concerns \textit{the method by which they assign weights to the gradients of the log likelihoods of tokens}. In GRPO, tokens receive weights based on their associated “importance weight” $\frac{ \pi_{\theta} (y_{i,t} | x, y_{i,<t}) }{ \pi_{\theta_\text{old}} (y_{i,t} | x, y_{i,<t})}$. These weights are not uniform; for cases where $\widehat{A}_{i} > 0$, they may range within $(0, 1+\varepsilon]$, while for $\widehat{A}_{i} < 0$, the weights can span $[1-\varepsilon, +\infty)$. Because these weights are not consistent across tokens, their influence can build up and potentially introduce instability as training continues. On the other hand, GSPO applies identical weighting to every token in a response, thus avoiding the potential for instability found in GRPO.

\subsection{GSPO-token: A Token-level Objective Variant}

In scenarios like multi-turn RL, we may desire a finer-grained advantage adjustment than the sequence level. To this end, we introduce a token-level objective variant of GSPO, namely \textbf{GSPO-token}, to allow token-wise advantage customization:

\begin{align}
\mathcal{J}_\text{GSPO-token}(\theta)
=
\mathbb{E}_{ x \sim \mathcal{D},\, \{y_i\}_{i=1}^G \sim \pi_{\theta_\text{old}}( \cdot | x) }
\left[ 
\frac{1}{G} \sum_{i=1}^{G} \frac{1}{|y_i|} \sum_{t=1}^{|y_i|} 
\min \left( s_{i,t}(\theta) \widehat{A}_{i,t},  \, \mathrm{clip} \left( s_{i,t}(\theta), 1 - {\varepsilon}, 1 + {\varepsilon} \right) \widehat{A}_{i,t} \right) 
\right],
\label{equ:gspo-token}
\end{align}

where

\begin{align}
s_{i,t}(\theta) 
= \mathrm{sg} \left[ s_{i}(\theta) \right]  \cdot \frac{ \pi_{\theta} (y_{i,t} | x, y_{i,<t}) }{ \mathrm{sg} \left[ \pi_{\theta} (y_{i,t} | x, y_{i,<t}) \right] },
\end{align}

and $\mathrm{sg}[\cdot]$ denotes only taking the numerical value but stopping the gradient, corresponding to the \texttt{detach} operation in PyTorch. The gradient of GSPO-token can be derived as:

\begin{align}
\nabla_{\theta} \mathcal{J}_\text{GSPO-token}(\theta)
=&\ 
\nabla_{\theta} \mathbb{E}_{ x \sim \mathcal{D},\, \{y_i\}_{i=1}^G \sim \pi_{\theta_\text{old}}( \cdot | x) }
\left[ 
\frac{1}{G} \sum_{i=1}^{G}
\frac{1}{|y_i|} \sum_{t=1}^{|y_i|} 
s_{i,t}(\theta) \widehat{A}_{i,t}
\right] \\
=&\ 
\mathbb{E}_{ x \sim \mathcal{D},\, \{y_i\}_{i=1}^G \sim \pi_{\theta_\text{old}}( \cdot | x) }
\left[ 
\frac{1}{G} \sum_{i=1}^{G} s_i (\theta) \cdot \frac{1}{|y_i|} \sum_{t=1}^{|y_i|} 
\widehat{A}_{i,t} \frac{ \nabla_{\theta} \pi_{\theta} (y_{i,t} | x, y_{i,<t}) }{ \pi_{\theta} (y_{i,t} | x, y_{i,<t}) }
\right] \\
=&\ 
\mathbb{E}_{ x \sim \mathcal{D},\, \{y_i\}_{i=1}^G \sim \pi_{\theta_\text{old}}( \cdot | x) }
\left[ 
\frac{1}{G} \sum_{i=1}^{G}  \left( \frac{ \pi_{\theta} (y_i | x) }{ \pi_{\theta_\text{old}} (y_i | x)} \right)^{\frac{1}{|y_i|}}
\cdot \frac{1}{|y_i|} \sum_{t=1}^{|y_i|}
\widehat{A}_{i,t} \nabla_{\theta} \log \pi_{\theta} (y_{i,t} | x, y_{i,<t})  
\right].
\label{equ:gspo-token_grad}
\end{align}

Observe that the expression $\frac{ \pi_{\theta} (y_{i,t} | x, y_{i,<t}) }{ \mathrm{sg} \left[ \pi_{\theta} (y_{i,t} | x, y_{i,<t}) \right] }$ evaluates to 1, which means that $s_{i,t}(\theta)$ coincides numerically with $s_{i}(\theta)$. By examining Equation~\eqref{equ:gspo} alongside Equation~\eqref{equ:gspo-token}, as well as Equation~\eqref{equ:gspo_grad} with Equation~\eqref{equ:gspo-token_grad}, it can be seen that GSPO-token and GSPO are equivalent in terms of numerical values for the optimization objective, clipping criteria, and analytical gradient calculations when all tokens in a given response $y_i$ share a uniform advantage (that is, $\widehat{A}_{i,t} = \widehat{A}_{i}$ for every token). However, GSPO-token allows for greater adaptability by permitting token-level variations in advantage values.

\section{Experiments and Discussion}

\subsection{Empirical Results}

We conduct experiments using a cold-start model, fine-tuned from Qwen3-30B-A3B-Base, and present both the progression of training rewards and performance trends on several benchmarks: AIME'24 (mean Pass@1 across 32 samples), LiveCodeBench (202410-202502, mean Pass@1 from 8 samples), and CodeForces (Elo Rating). During reinforcement learning (RL) training, we subdivide each rollout batch into four mini-batches for individual gradient update steps. For GSPO, the left and right clipping intervals in Equation~\eqref{equ:gspo} are assigned values of 3e-4 and 4e-4, respectively. For baseline comparison, we utilize GRPO, adjusting its clipping parameters in Equation~\eqref{equ:grpo} to 0.2 (left) and 0.27 (right); these have been meticulously optimized to facilitate an equitable comparison. It is important to highlight that, for stable convergence in MoE RL, GRPO requires the Routing Replay training approach—which will be further addressed in \S~\ref{subsec:routing_replay}. In contrast, \textit{GSPO eliminates the necessity for this additional strategy}.

Figure~\ref{fig:results} illustrates that training with GSPO exhibits consistent stability. Notably, \textbf{GSPO enables ongoing gains in performance by scaling up training computation, frequently refreshing the query set, and allowing for longer generation sequences}. Additionally, GSPO outperforms GRPO in terms of training efficiency, attaining higher training accuracy and improved benchmark results while using equivalent amounts of computation and query resources. Furthermore, we have effectively implemented GSPO in the RL training process of the state-of-the-art Qwen3 models, providing strong evidence for GSPO’s capability to harness RL scaling benefits in large language model development.

\subsection{Curious Observation on Clipping Fractions}

An important difference between GSPO and GRPO lies in their respective clipping strategies: GSPO applies clipping at the sequence (response) level, whereas GRPO clips on a per-token basis. Notably, as depicted in Figure~\ref{fig:clipping}, the proportion of tokens clipped under GSPO surpasses that of GRPO by approximately two orders of magnitude—a pattern that persists even when modifying the clipping thresholds. Nevertheless, even though GSPO clips a substantially larger number of tokens and, as a result, relies on fewer tokens for training or gradient estimation, it still outperforms GRPO in terms of training efficiency. This somewhat surprising outcome—whereby greater proportions of token clipping enhance, rather than hinder, training effectiveness—suggests that the token-level gradient approximations used by GRPO are fundamentally noisy and ill-suited for sample efficiency. Conversely, the sequence-level gradients obtained through GSPO furnish a more stable and potent signal for learning.

\subsection{Benefit of GSPO for MoE Training}
\label{subsec:routing_replay}

\paragraph{Background}
Relative to reinforcement learning (RL) training with dense architectures, Mixture of Experts (MoE) models, due to their inherent sparsity in expert activation, present distinct challenges regarding training stability. Notably, our investigations reveal that when employing the GRPO algorithm, the characteristic instability in expert selection—termed \textit{expert-activation volatility}—can impede the convergence of RL training in MoE architectures. Concretely, following a single or several gradient update steps, the set of experts engaged for a repeated response may shift substantially. Taking the 48-layer Qwen3-30B-A3B-Base as an illustrative case, we observe that, after each RL gradient step and for the identical rollout instance, about 10\% of the experts chosen under the updated policy $\pi_{\theta}$ do not overlap with those under the former policy $\pi_{\theta_\text{old}}$. This pattern, which is even more evident as model depth increases in MoE models, exacerbates the variance of token-level importance weights $w_{i,t}(\theta) = \frac{ \pi_{\theta} (y_{i,t} | x, y_{i,<t}) }{ \pi_{\theta_\text{old}} (y_{i,t} | x,y_{i,<t})}$, leading to erratic and untrustworthy ratios, as elaborated in \S~\ref{sec:motivation} and~\ref{subsec:gradient}. As a result, these fluctuations seriously disrupt the RL optimization process, thwarting smooth and stable convergence.

\paragraph{Our Previous Approach} In our earlier work, we addressed this issue by implementing the \textbf{Routing Replay} training method. This involved recording the expert routes selected by $\pi_{\theta_\text{old}}$ and subsequently using these identical routing paths in $\pi_{\theta}$ when calculating the importance weights $w_{i,t}(\theta) = \frac{ \pi_{\theta} (y_{i,t} | x, y_{i,<t}) }{ \pi_{\theta_\text{old}} (y_{i,t} | x,y_{i,<t})}$. Through this procedure, both $\pi_{\theta} (y_{i,t} | x, y_{i,<t})$ and $\pi_{\theta_\text{old}} (y_{i,t} | x,y_{i,<t})$ process each token $y_{i,t}$ using the same sub-network. This consistency helps to stabilize the calculation of token-wise importance ratios and enables effective optimization by updating the same set of parameters throughout training steps. As illustrated in Figure~\ref{fig:routing_replay}, Routing Replay is pivotal for ensuring steady convergence of GRPO training in MoE frameworks.

\paragraph{Benefit of GSPO} While Routing Replay makes it possible for MoE models to achieve stable convergence with GRPO training, this method requires storing and sharing routing states, thereby increasing memory and communication costs and potentially restricting the model’s effective capacity. On the other hand, as demonstrated in Figure~\ref{fig:results}, GSPO obviates the need for Routing Replay entirely, allowing conventional computation of importance ratios $s_i(\theta)$ and supporting stable, normal convergence. The central idea underlying GSPO is that it relies exclusively on the sequence-level likelihood (i.e., $\pi_{\theta} (y_i | x)$), remaining unaffected by token-level variations (i.e., $\pi_{\theta} (y_{i,t} | x, y_{i,<t})$). Given that the MoE model consistently upholds its language modeling performance, significant variability in the sequence likelihood does not occur. Overall, GSPO directly addresses the instability in expert activation that plagues MoE training and removes the necessity for intricate strategies like Routing Replay. This leads to a streamlined and more reliable training regimen and permits the model to fully utilize its capacity without extrinsic limitations.

\subsection{Benefit of GSPO for RL Infrastructure}

Due to the differences in numerical precision between training engines such as Megatron and inference engines like SGLang and vLLM, it is common in practice to recompute the likelihoods of sampled outputs using the training engine for the old policy $\pi_{\theta_\text{old}}$. However, GSPO’s optimization relies on sequence-level likelihoods rather than token-level ones; sequence-level likelihoods are generally more robust to precision variation. As a result, GSPO enables the direct use of likelihoods produced by the inference engine during optimization, eliminating the need to perform recomputation with the training engine. This approach is particularly advantageous in settings such as partial rollout, multi-turn RL, and frameworks where training and inference processes are separated.

\section{Conclusion}

We introduce Group Sequence Policy Optimization (GSPO), a reinforcement learning algorithm specifically designed for large language model training. GSPO leverages importance sampling principles by constructing importance ratios from sequence likelihoods, and executes operations such as sequence-level clipping, rewarding, and optimization. Compared with GRPO, GSPO offers enhanced stability during training, increased efficiency, and improved performance, showing particular strength in large-scale RL training scenarios for MoE models. This robust performance underpins the significant advancements seen in the latest Qwen3 models. As a scalable method, GSPO forms the algorithmic backbone for ongoing RL scaling efforts, which we anticipate will drive further foundational progress in intelligence.

\end{document}